En 1935, Albert Einstein, Boris Podolsky y Nathan Rosen (EPR) hicieron manifiesto de su postura sobre la teoría de la Mecánica Cuántica. Ellos sugieren que la descripción cuántico-mecánica de la realidad dada por las funciones de onda no es completadejando abierta la pregunta sobre si existe o no tal descripción que hagan de la mecánica cuántica una teoría completa. \cite{EPR} Este artículo, a pesar de sí tener una parte matemática, llega a rozar lo "filosófico" sobre lo que implica y lo que no implica la teoría de la Mecánica Cuántica, introduciendo un concepto (no exahustivo) sobre lo que implica que algo sea "real".

En 1964, John Bell, con estos dos desarrollos (más algunos varios más) en el primer volumen de la revista Pyhsics, formula lo que ahora se conocen como la desigualdad de Bell. Es la formulación matemática que se debe de satisfacer si un sistema es local y real. 

Años más tarde, numerosos autores deducen variantes de esta desigualdad, sin embargo, la suposición de un realismo local no se presenta como un paso lógico específico para derivar las desigualdades. Es por ello que en este trabajo, no sin antes enunciar una forma de la "paradoja" EPR, se retomará el caso de un estado entrelazado de dos partículas y se hará una deducción de la desigualdad original de Bell sin pasar por una teoría de variables ocultas de manera explícita, y apoyándose de un método gráfico realizado por Krzysztof Rebilas \cite{Rebilas} para desarrollarlo de una manera más intuitiva. Finalmente, se abordará qué implicaciones tiene que no se cumpla esa desigualdad que planteó Bell bajo estas hipótesis.